% Prompts usados na geração dos testbenches (kit_prompts_testbench, colado pelo autor em 2026-11-09).
% O texto é o mesmo enviado aos LLMs; só a indentação do código e as quebras das linhas longas
% foram ajustadas para caber na página. Sem tabela, a pedido do autor: a composição vai em lista.

\section{\textit{Prompts} utilizados}
\label{ap:prompts}

Este apêndice apresenta os seis \textit{prompts} usados na geração dos \textit{testbenches}. Como eles compartilham trechos, os trechos comuns são mostrados uma única vez, depois da lista abaixo, que indica a ordem em que cada \textit{prompt} foi montado. Em todos eles, \texttt{[MÓDULO]} foi substituído pelo código RTL completo do módulo avaliado. A indentação do código e as quebras das linhas longas foram ajustadas para caber na página.

\begin{itemize}
  \item \textbf{\textit{Zero-shot}}: ``Crie um testbench para o módulo RTL abaixo.'' seguido de [Módulo e requisitos].
  \item \textbf{\textit{Few-shot}}: [Exemplos], ``Agora crie um testbench para o módulo RTL abaixo.'' e [Módulo e requisitos].
  \item \textbf{Especialista \textit{Zero-shot}}: ``Aja como um engenheiro de verificação de hardware. Crie um testbench para o módulo RTL abaixo.'' seguido de [Módulo e requisitos].
  \item \textbf{Especialista \textit{Few-shot}}: ``Aja como um engenheiro de verificação de hardware.'', [Exemplos], ``Agora crie um testbench para o módulo RTL abaixo.'' e [Módulo e requisitos].
  \item \textbf{\textit{Planning Zero-shot}}: [Planejamento] seguido de [Módulo e requisitos].
  \item \textbf{\textit{Planning Few-shot}}: [Exemplos], [Planejamento] e [Módulo e requisitos].
\end{itemize}

\noindent\textbf{[Exemplos]}, usado nas três estratégias \textit{few-shot}:

{\footnotesize
\begin{verbatim}
A seguir estão exemplos de módulos RTL e testbenches corretos. Use-os apenas como referência
de estrutura, estilo e forma de verificação.

Exemplo 1:
module comparador4 (
    input  [3:0] a,
    input  [3:0] b,
    output       maior,
    output       igual,
    output       menor
);
    assign maior = (a > b);
    assign igual = (a == b);
    assign menor = (a < b);
endmodule

module tb;
    reg  [3:0] a, b;
    wire maior, igual, menor;
    integer erros = 0;

    comparador4 dut (.a(a), .b(b), .maior(maior), .igual(igual), .menor(menor));

    initial begin
        a = 4'd5; b = 4'd3; #1;
        if ({maior, igual, menor} !== 3'b100) erros = erros + 1;
        a = 4'd2; b = 4'd9; #1;
        if ({maior, igual, menor} !== 3'b001) erros = erros + 1;
        a = 4'd7; b = 4'd7; #1;
        if ({maior, igual, menor} !== 3'b010) erros = erros + 1;
        if (erros != 0) $fatal(1, "comparador: %0d erros", erros);
        $display("[PASSOU] comparador4 ok");
        $finish;
    end
endmodule

Exemplo 2:
module shift_reg4 (
    input        clk,
    input        rst,
    input        serial_in,
    output [3:0] q
);
    reg [3:0] reg_q;

    always @(posedge clk) begin
        if (rst)
            reg_q <= 4'd0;
        else
            reg_q <= {reg_q[2:0], serial_in};
    end

    assign q = reg_q;
endmodule

module tb;
    reg clk = 0, rst, serial_in;
    wire [3:0] q;
    integer erros = 0;

    shift_reg4 dut (.clk(clk), .rst(rst), .serial_in(serial_in), .q(q));

    always #5 clk = ~clk;
    task step; begin @(posedge clk); #1; end endtask

    initial begin
        rst = 1; serial_in = 0; step;
        if (q !== 4'd0) erros = erros + 1;
        rst = 0;
        serial_in = 1; step;
        if (q !== 4'b0001) erros = erros + 1;
        serial_in = 0; step;
        if (q !== 4'b0010) erros = erros + 1;
        serial_in = 1; step;
        if (q !== 4'b0101) erros = erros + 1;
        serial_in = 1; step;
        if (q !== 4'b1011) erros = erros + 1;
        if (erros != 0) $fatal(1, "shift: %0d erros", erros);
        $display("[PASSOU] shift_reg4 ok");
        $finish;
    end
endmodule
\end{verbatim}
}

\noindent\textbf{[Planejamento]}, usado nas duas estratégias \textit{Planning}:

{\footnotesize
\begin{verbatim}
Antes de gerar o testbench, analise o módulo RTL abaixo. Primeiro, identifique as entradas e
saídas do módulo, verifique se ele é combinacional ou sequencial, descreva o comportamento
esperado, identifique estados, operações, ramos, sinais de controle, casos de borda e
cenários de teste.
Depois da análise, gere o testbench final.
\end{verbatim}
}

\noindent\textbf{[Módulo e requisitos]}, usado nos seis \textit{prompts}:

{\footnotesize
\begin{verbatim}
Módulo RTL:
[MÓDULO]

Requisitos do testbench:
- O módulo do testbench deve se chamar tb.
- Instancie o módulo sob teste e compare as saídas com os valores esperados usando !==.
- Em caso de erro, use $fatal(1, ...) para encerrar com código de erro.
- Cada $display deve ter a string em uma única linha.
- Termine com $finish.
- Deve compilar no Verilator (código sintetizável, sem arquivos externos).
\end{verbatim}
}
